problem:  
Let \( X \) be a smooth Fano manifold, and suppose \( \mathcal{X} \to \Delta \) is a flat, \(\mathbb{Q}\)-Gorenstein smoothing whose central fiber \(X_0\) is a klt \(\mathbb{Q}\)-Fano variety that is **not known or assumed to be \(K\)-polystable**. For \(t\neq 0\), let \(\omega_t\) denote the unique Kähler–Einstein metric on \(X_t\), normalized to satisfy \(\mathrm{Ric}(\omega_t)=\omega_t\).  

Consider the family of pointed metric spaces \((X_t, \omega_t)\) as \(t\to 0\).  

**Question:**  
Does there exist a canonical limiting metric space \((Z, d_Z)\) that is uniquely determined (independent of the smoothing family \(\mathcal{X}\to \Delta\)) by the birational class of \(X_0\)? In particular, if two non-isomorphic \(\mathbb{Q}\)-Gorenstein smoothings of the same singular \(\mathbb{Q}\)-Fano variety \(X_0\) both admit smooth fibers with Kähler–Einstein metrics, must their Gromov–Hausdorff limits coincide up to isometry?  

Equivalently, is the Gromov–Hausdorff compactification of Kähler–Einstein Fano manifolds functorial with respect to birational maps among their degenerations, even beyond the realm of \(K\)-polystability?  

Why is it a "good" problem:  
This problem isolates one subtle frontier beyond existing Chen–Donaldson–Sun theory. The known correspondence between algebraic and metric compactifications of Kähler–Einstein Fano manifolds crucially depends on \(K\)-polystability. In contrast, this question probes whether the uniqueness and canonical nature of metric limits persist **without** the \(K\)-polystability assumption, i.e., whether metric compactifications encode a deeper birational functoriality still invisible to current algebro-geometric frameworks.  

It bridges three perspectives—birational geometry, metric convergence, and the stability theory of Fano manifolds—asking whether purely analytic limits “remember” algebraic birational structures even when stability fails. The question is precise, unproven, and conceptually simple yet profound: it challenges the sufficiency of existing moduli compactification theories and suggests a possible intrinsic metric definition of canonical degenerations independent of \(K\)-stability machinery. Its resolution would reshape our understanding of how analytic and algebro-geometric moduli interact at their boundaries.